% !TeX root = entrypoint.tex
\documentclass[oneside,notitlepage,12pt]{article}
\pagestyle{plain}

%\usepackage{polish}
\usepackage[british,UKenglish,USenglish,american]{babel}
\usepackage[utf8]{inputenc}

\usepackage{amssymb}
\usepackage[leqno]{amsmath}
\usepackage{amsfonts}
\usepackage{amsopn}
\usepackage{amstext}
\usepackage{amsthm}
\usepackage[textsize=small]{todonotes}
%\usepackage[disable]{todonotes}
\setuptodonotes{color=blue!30}
\usepackage{enumitem}

\usepackage{tikz}
\usetikzlibrary{cd}

\usepackage{verbatim}
\usepackage[colorlinks, backref, citecolor=blue]{hyperref}
\usepackage[numbers]{natbib}
\usepackage{makeidx}

\newcommand{\define}[2]{{\em #1}\index{#2}}


% Packages for special symbols and ornaments:
\usepackage{calrsfs}
\usepackage{fourier-orns}
%\usepackage{hieroglf} \newcommand{\hiero}{\textpmhg}
\usepackage{clock} %\ClockFrametrue\ClockStyle2
\usepackage[alpine, weather]{ifsym}
\usepackage{tabularx} % extra features for tabular environment
\usepackage{graphicx} % takes care of graphic including machinery
\usepackage{blindtext}
\usepackage{soul}


%PAGE SETUP
\textheight=22cm
\textwidth=15cm
%\hoffset=-1cm
\hoffset=1cm
\voffset=-2cm
%\parindent=16pt
\setlength{\marginparwidth}{4cm}
\reversemarginpar

% Notes in margins and switching them on and off (if needed)
\newcommand{\almarginpar}[1]{\leavevmode\marginpar[\raggedleft\small\em{#1}]{}}
%\newcommand{\almarginpar}[1]{}

\providecommand{\ar}{\arrow}
\newcommand{\Em}{\mathbb M}


\frenchspacing

\providecommand{\cal}{\mathcal}
\renewcommand{\Bbb}{\mathbb}
\renewcommand{\frak}{\mathfrak}
\newenvironment{pf}{\begin{proof}}{\end{proof}}

%%%%%%%%%%%%%%%%%%%%
% Standard commands
%%%%%%%%%%%%%%%%%%%%

%%%%%%%%%%%%%%%%%%%%
% Calligraphic and bold letters. 
%%%%%%%%%%%%%%%%%%%%
\newcommand{\Aa}{{\Bbb{A}}}
\newcommand{\Aaa}{{\cal{A}}}
\newcommand{\Bee}{{\cal{B}}}
\newcommand{\Cee}{{\cal{C}}}
\newcommand{\Dee}{{\cal{D}}}
\newcommand{\Ef}{{\cal{F}}}
\newcommand{\Gee}{{\cal{G}}}
\newcommand{\Haa}{{\cal{H}}}
\newcommand{\Ai}{\cal{I}}
\newcommand{\Kay}{{\cal{K}}}
\newcommand{\El}{{\cal{L}}}
\newcommand{\Pee}{{\cal{P}}}
\newcommand{\See}{{\cal{S}}}
\newcommand{\Tau}{{\cal{T}}}
\newcommand{\Yu}{{\cal{U}}}
\newcommand{\Vee}{{\cal{V}}}
\newcommand{\Wu}{{\cal{W}}}
\newcommand{\Zee}{{\Bbb{Z}}}
\newcommand{\Emm}{{\frak{M}}}
\newcommand{\Be}{{\Bbb{B}}}
\newcommand{\Nat}{{\Bbb{N}}}
\newcommand{\Z}{{\mathbb Z}} % The integers
\newcommand{\Qyu}{{\Bbb{Q}}}
\newcommand{\Err}{{\Bbb{R}}}


%%%%%%%%%%%%%%%%%%%%
% Shortcuts for some Greek letters. 
%%%%%%%%%%%%%%%%%%%%
\newcommand{\lam}{{\lambda}}
\newcommand{\al}{\alpha}
\newcommand{\Gam}{\Gamma}
\newcommand{\alG}{\al\in\Gamma}
\newcommand{\sig}{\sigma}
\newcommand{\eps}{\varepsilon}
\renewcommand{\phi}{\varphi}
\renewcommand{\rho}{\varrho}

%%%%%%%%%%%%%%%%%%%%
% Basic commands. 
%%%%%%%%%%%%%%%%%%%%
\newcommand{\rest}{\restriction}
\newcommand{\unii}{\mathbb I}
\newcommand{\ntr}{{n\in\omega}}
\newcommand{\Ntr}{n\in{\Bbb{N}}}
\newcommand{\loe}{\leq}
\newcommand{\goe}{\geq}
\newcommand{\notloe}{\nleq}
\newcommand{\notgoe}{\ngeq}
\newcommand{\subs}{\subseteq}
\newcommand{\sups}{\supseteq}
\newcommand{\nnempty}{\ne\emptyset}
\newcommand{\argum}{\:\cdot\:}
\newcommand{\ovr}{\overline}
\renewcommand{\iff}{\Longleftrightarrow}


%%%%%%%%%%%%%%%%%%%%
% Topology. 
%%%%%%%%%%%%%%%%%%%%
\newcommand{\Cl}[1]{\overline{#1}}
\newcommand{\cl}{\operatorname{cl}}
\newcommand{\Int}{\operatorname{int}}
\newcommand{\w}{\operatorname{w}}
\newcommand{\dens}{\operatorname{dens}}
\newcommand{\eXp}{\operatorname{exp}}
\newcommand{\diam}{\operatorname{diam}}
\newcommand{\dist}{\operatorname{dist}}
\newcommand{\nbd}{\operatorname{nbd}}



%%%%%%%%%%%%%%%%%%%%
% Convexity. 
%%%%%%%%%%%%%%%%%%%%
\newcommand{\conv}{\operatorname{conv}}
\newcommand{\co}{\operatorname{co}}
\newcommand{\clco}{\operatorname{clco}}
\newcommand{\m}{\operatorname{m}}
\newcommand{\med}{\operatorname{med}}

\newcommand{\G}{{\mathbb G}}

%%%%%%%%%%%%%%%%%%%%
% Miscellanous. 
%%%%%%%%%%%%%%%%%%%%
\newcommand{\id}[1]{{\operatorname{i\!d}_{#1}}} % identity morphism
\newcommand{\cf}{\operatorname{cf}}
\newcommand{\dom}{\operatorname{dom}}
\newcommand{\cod}{\operatorname{cod}}
\newcommand{\Dom}{\operatorname{Dom}}
\newcommand{\rng}{\operatorname{rng}}
\newcommand{\suppt}{\operatorname{suppt}}
\newcommand{\defi}{\stackrel{\rm df}=}
\newcommand{\pr}{\operatorname{pr}}
\newcommand{\liminv}{\varprojlim}

\newcommand{\symd}{\div} % <--- Symmetrical difference

\newcommand{\oraz}{\qquad\text{and}\qquad}

%%%%%%%%%%%%%%%%%%%%
% Some forcing commands. 
%%%%%%%%%%%%%%%%%%%%
\newcommand{\proves}{\vdash}
\newcommand{\forces}{\Vdash}
\newcommand{\Con}{\operatorname{Con}}
\newcommand{\Card}{{\frak{Card}}}
\newcommand{\Func}{\operatorname{Func}}
\newcommand{\rank}{\operatorname{rank}}
\newcommand{\poset}{{\Bbb{P}}}
\newcommand{\Forall}{\forall\;}
\newcommand{\Exists}{\exists\;}
\newcommand{\h}{\widehat}
\newcommand{\val}{\operatorname{val}}
\newcommand{\comp}{\parallel}
\newcommand{\incomp}{\perp}
\newcommand{\Es}{{\cal{S}}}
\newcommand{\meet}{\wedge}
\newcommand{\Meet}{\prod}
\newcommand{\join}{\vee}
\renewcommand{\Join}{\sum}
\newcommand{\Land}{\;\&\;}


%%%%%%%%%%%%%%%%%%%%
% Trees. 
%%%%%%%%%%%%%%%%%%%%
\newcommand{\Lev}{\operatorname{Lev}}
\newcommand{\lev}[2]{{\ell_{#1}#2}}
\newcommand{\Ht}{\operatorname{ht}}
\newcommand{\length}{\operatorname{length}}
\newcommand{\bd}{\partial}

\newcommand{\concat}{{}^\smallfrown}
\newcommand{\Reg}{{\frak{Reg}}}
\newcommand{\Ord}{{\frak{Ord}}}
\newcommand{\by}[1]{/_{#1}}
\newcommand{\club}{\operatorname{club}}
\newcommand{\proj}{\operatorname{proj}}
\newcommand{\otp}{\operatorname{otp}}
\newcommand{\Fin}{\operatorname{fin}}

%%%%%%%%%%%%%%%%%%%%
% Theorems and Propositions. 
%%%%%%%%%%%%%%%%%%%%
\newtheorem{tw}{Theorem}[section]
\newtheorem{wn}[tw]{Corollary}
\newtheorem{lm}[tw]{Lemma}
\newtheorem{prop}[tw]{Proposition}
\newtheorem{claim}[tw]{Claim}
\theoremstyle{definition}
\newtheorem{df}[tw]{Definition}
\newtheorem{ex}[tw]{Example}
\newtheorem{pyt}[tw]{Question}
\newtheorem{question}[tw]{Question}
\theoremstyle{remark}
\newtheorem{uwgi}{Remark}
\renewcommand{\theuwgi}{}

% Cantor's stuff:
\newcommand{\Cantor}{2^\omega}

\hypersetup{
	colorlinks=true,       % false: boxed links; true: colored links
	linkcolor=blue,        % color of internal links
	citecolor=blue,        % color of links to bibliography
	filecolor=magenta,     % color of file links
	urlcolor=blue         
}

%++++++++++++++++++++++++++++++++++++++++


\begin{document}

	% % \begin{titlepage}
    % \begin{center}
    %     \vspace*{1cm}
            
    %     \Huge
    %     \textbf{Exploration of Variational Quantum Algorithm}
            
    %     \vspace{0.5cm}
    %     \LARGE
    %     \title{I will pick a title later}
        
            
    %     \vspace{1.5cm}
            
    %     \textbf{Thanh Nguyen and Jacob Cybulski}
            
    %     \vfill
            
    %     A thesis presented for the degree of\\
    %     Bachelor of IT Honours
            
    %     \vspace{0.8cm}
            
    %     \includegraphics[width=0.4\textwidth]{src/CoverPage/DeakinUniversityLogo.jpg}
        
            
    %     \Large
    %     % Faculty of Science, Engineering and Built Environment\\
    %     Deakin University\\
    %     \date{\today}
            
    % \end{center}
% \end{titlepage}  
\thispagestyle{empty}
\begin{titlepage}
    \includegraphics[width=0.25\textwidth]{src/CoverPage/Deakin_Logo.jpeg}
        \begin{center}
        \vspace*{4cm}
        \todo{May change later on}
        {\LARGE Investigation of Barren Plateaus in Quantum Neural Network development} %%Replace this with the Title of your research
        \vspace{3cm}
            \begin{large}   
    
        
            \bf Submitted as Honours Dissertation in SIT723/SIT724
            \vspace{1cm}
        
            \bf \today \\
            T1-2022        
        
            \vspace{3cm}
            \textbf{Thanh Nguyen}\\
            STUDENT ID 218583133 \\
            COURSE - Bachelor of IT Honours (S470)
            \vfill

            \bf \normalsize Supervised by: Prof, Jacob Cybulski\\
       
        \end{large}  
   \end{center}
\end{titlepage}


% \begin{abstract}
Quantum Computing is a new and exciting area of Science, which intersects Physics, Mathematics and Computer Science. Quantum algorithms are known to solve problems, which until recently have been considered unsolvable classically due to their inherent complexity. Typically a high-level quantum algorithm in Python generates a quantum circuit, which can be executed on a real quantum machine or a simulator. Unfortunately, quantum circuits hard-code their data, so that an application of an algorithm to new data leads to a new circuit. This means that circuits are static and cannot be trained or optimised directly. However, there exist hybrid classical-quantum techniques that can produce variational (or parametrised) circuits, which act as templates for circuit generation. Quantum machine learning algorithms take advantage of variational circuits to implement quantum-alternatives to many machine learning algorithms. In this way, quantum solutions can be both highly efficient and data rich.

This project aims to investigate a range of variational quantum algorithms, i.e. hybrid classical-quantum algorithms capable of manipulating parametrised quantum circuits.
\\\\
{\bf Keywords:} Quantum Computing, Quantum Machine Learning, Variational Quantum Algorithms, Python Programming, Qiskit.

\end{abstract}

% \tableofcontents

\section{Project Description}

\subsection{Background}

\todo{you really need to fix your English throughout}Regardless of \todo[prepend]{vast?}\underline{far-fetched} power of supercomputers, there are problems that \todo{just "remain"}\underline{are remaining} unsolved, as their complexity scales too quickly even for some top-of-the-line computing systems to handle. 
We are in needed to research and develop new tools with enough processing power to process those tasks, which can be taken forms of molecular behaviors analyzing in chemistry, factoring large integer for cryptography, or machine learning.

Quantum computers were developed based on Quantum mechanics to deliver a greater processing scaling compare to Classical computers.
It can be considered as the relic of the collaboration between modern physics, engineering, and computer science.
To demonstrate the advantages of Quantum computers, Google claimed "Quantum supremacy" in 2019, by solving a problem that take about 10000 years for supercomputers, in 200 seconds \cite{hsuGoogleQuantumTech2019}.
New records for this Quantum race are established over and over, such as Jiuzhang Quantum computer from China with 56 Qbits configuration, reported to have 10 billion times in performance compare to Google's \cite{zhongQuantumComputationalAdvantage2020}, or the most recent 127 Qbits system "Eagle" from IBM \cite{chow2021ibm}.

Both Classical and Quantum computer suffer the same fundamental issue of data corruption.
For Classical computers, the cause of this fault would be magnetism, electricity and radiations.
However, in the case of Quantum computers, there are even more things to concern, not only from magnetism or electricity and radiations, they require separation from heat, light, vibration, etc. In other words, an absolute isolation from the outside world.

Such close systems are of course only exist in fiction. 
All Quantum computers must be designed to be controlled by interaction with the outside world at some extent, or else we cannot input or retrieve data.
It also means there are possibility that some disturbances that we have no control of can interrupt the system, or the system itself interacting with the outside world.
These interruptions, called "decoherence", are the major factor that destroy the Quantum evolution needed for large scale algorithms on Quantum computers.

A way to counter decoherence is using Quantum Error Correction codes \cite{lidar2013quantum} to protect quantum information, by using multiple physical Qbits to represent a logical Qbit, such as the 9-Qbit Shor code \cite{shor1995scheme}.
Considering the high cost of fault-tolerant Quantum computers and their maintenance, Qbits must be used as efficient as possible. 
The limited in amount of Qbits per processor and circuit depth due to noise at this time also narrowing down on the scale of algorithms that can be executed.

The availability of those fault-tolerant Quantum computer is not likely to happen in the near future. 
While waiting for those devices, scientists had discovered a new approach for the current early and noisy Qbit registers.
Variational Quantum algorithms (VQAs) \cite{cerezo2021variational} attempted to address this constrains, by made used of machine learning methods to train a parameterised quantum circuit.
\todo{the preceding discussion does not address the issue of Variational Quantum Algorithms, the fact that QNN is an example of VQA, the reason for investigating these, etc. }
\begin{center}
    \vspace{10pt}
    Add more later
    \vspace{10pt} 
\end{center}
Perhaps, something like this should be in this section: You need more focus, do the following - (1) open with a general statement about quantum computing and their benefits as compared with classical computers, (2) explain how quantum algorithms are being developed and the fact that they are "static", (3) discuss that an alternative involved hybrid solutions, such as VQA and give some examples of VQAs, including QNNs, (4) briefly discuss/mention different types of QNNs, their advantages and disadvantages, (5) set a claim that this project undertakes a study aims to compare different approaches to implementing QNNs and contrasting them with classical NNs, (6) finally mention that the research will focus on a classification problem (e.g. of images).


\todo{May need to narrow down the scope, we only have 8 months, most likely 6 months for writing and another 2 for editing and submitting}
\subsection{The problem}
Here are some topics (need to pick one):
\begin{enumerate}
    \item Implementation of Binary image classification with objects of distinct classes $A$ and $B$. For example: Cats and Dogs.
    \item Feature maps of high-order in relation with training capability as addressed in \cite{abbasPowerQuantumNeural2021}.
    \item The Barren Plateaus in QNN circuit training and how to avoid/deal with them
\end{enumerate}
\todo{2 and 3 are interesting, but you ideally need to formulate the problem into a question, e.g. "How..."}
Topics other than QNN:

Solving Max2Xor / knapsack / max-cut with QAOA.

\subsection{Objectives}
\begin{enumerate}
    \item Pick a name (Due this weekend),
    \item Run through Classical Machine Learning (briefly),
    \item Hypothesis for QNN: From Classical Machine Learning to VQA (model, feature map, anzats, cost function, optimizers, etc.), a model that maps objects of distinct classes to binary output like: $A \to 0$ and $B \to 1$.
    \item Practice QNN: Image classification in python (Qiskit), maybe classifying animals, cats and mouses, using "Supervised learning" or any other settings. This step test the model if not work then need to alter the hypothesis.
    \item Conclude the research with summary, motivation and applications.
\end{enumerate}

Hypothesis for Feature maps and QCNN circuit training: A feature is an individual measurable property (characteristic in short) of a class of objects. A \textit{Feature map} is related to dimensionality reduction, it re-represent data in different space



\subsection{Success Criteria}
\begin{enumerate}
    \item This thesis/research focuses on exploring VQAs, more specifically Quantum Neural Network (QNN) and classification problem;
    \item A demonstration of a classifier that runs in a quantum computer.
\end{enumerate}



\subsection{Motivation}
\todo{We may need to work on these more, e.g. if you work on feature maps the motivation may be to provide a set of guidelines for the creation of feature maps suitable for QNNs and a given data set, if it is about barren plateaus, perhaps it is about finding a QNN architecture which eliminates barren plateaus}
\begin{enumerate}
    \item Learning ML with Quantum computer;
    \item This research also provides an approach to Quantum algorithms for ML engineers, and Computer science;
    \item Some useful applications with QNN:
    \begin{enumerate}
        \item Can work with early state of quantum computers (limited Qbit number and circuit depth);
        \item Error Correction is not necessary.
    \end{enumerate}
    \item Prove that VQAs (and QNN) can still be a promising field even for fault-tolerant Quantum computers in the future.
    \item \st{For my Bachelor of IT Hons degree ofc.}
\end{enumerate}


\bibliographystyle{jcabbrv} % Sorted and "note" fields removed
\bibliography{src/References/zoteroReferences.bib}

\end{document}
